% III. Програмна реализация.
\section{Програмна реализация}
\label{sec:implementation}

Реализацията е на C++20~\cite{iso_cpp20}: съпрограмите носят асинхронния модел, без да
раздробяват логиката на протокола между обратни извиквания; \code{concept} изразява общото
между двата кодека без виртуални функции; \code{std::span} и \code{std::string\_view} са
носителят на декодирането без копие. Сглобяването е с CMake и проектът няма нито една
задължителна външна зависимост: на чиста машина стигат компилатор за C++20 и CMake. Това има
цена, платена на три места. Тестовата
рамка е собствена, 170 реда в \code{tests/check.hpp}, свързана с CTest. Измерването е със
собствен часовник върху \code{std::chrono::steady\_clock}. Хеш функциите, от които зависят
двата обмена за удостоверяване, са в \code{src/crypto.cpp}: MD5, SHA-1 и SHA-256~\cite{fips180},
HMAC и PBKDF2, всяка проверена срещу публикуваните контролни вектори, а не срещу собствения си
изход. Единствената незадължителна зависимост е \code{libpq}, търсена с
\code{find\_package(PostgreSQL QUIET)} и използвана само от измервателната програма.

\textbf{Непълно съобщение.} Транспортът може да върне произволна част от пакет, така че
кодекът трябва да може да спре и да продължи, без да губи състояние. Решението е кодекът да
няма състояние изобщо: \code{peek\_frame} гледа буфера и връща или цяло съобщение, или празна
стойност, без да консумира нищо, а консумирането е отделна стъпка след обработката. Тестът
\code{peek\_frame\_refuses\_to\_consume\_a\_partial\_message} подава всяко възможно начало на
едно съобщение, от нула до пълната дължина, и изисква празна стойност за всяко. Четенето на
съобщение е изчакваем обект и \code{await\_ready()} вика \code{peek\_frame}, тоест в съпрограма
се влиза веднъж на четене от сокета, а не веднъж на съобщение (Раздел~\ref{sec:defects}).

\textbf{Грешка от сървъра в средата на резултатен поток.} Грешка, съобщена по правилата на
протокола, задължава клиента да изчерпи обмена до синхронизиращата точка, преди връзката да
бъде използвана отново. Реализацията записва грешката, но продължава да чете до
\code{ReadyForQuery} и чак тогава хвърля изключение. Затова връзка, върху която е изпълнено
грешно заявление, остава годна: интеграционният тест
\code{pg\_live\_server\_error\_leaves\_the\_connection\_usable} след грешка с код 42P01
изпълнява ново заявление по същата връзка.

\textbf{Заемане на връзка от пула.} Заемането се състезава с изтичане на срок. Съпрограмата
не може да изчаква едновременно свободна връзка и таймер, затова едното от двете идва като
обратно извикване, което урежда другото. Общото състояние се държи през \code{shared\_ptr}, но
извикването по срок държи \code{std::weak\_ptr}, защото не се отменя при нормално даване на
връзка и доживява до срока си (Раздел~\ref{sec:defects}).
Освобождаването предава връзката пряко на чакащия. Първата реализация я връщаше в списъка на
свободните и събуждаше чакащия, което изглежда еквивалентно. Не е: съпрограмата, която
освобождава връзката, продължава да се изпълнява преди събудения чакащ и си взема същата
връзка обратно, така че единият работник завършва целия си дял, преди другият да получи
изобщо една връзка. Поведението е записано като тест, а не като число:
\code{pool\_hands\_a\_released\_connection\_straight\_to\_the\_waiter} пуска два работника по
четири заемания при пул с размер едно и изисква ред на изпълнение \code{abababab}.

\textbf{Обезсилване на кеша.} Изместване от кеша връща изместения запис, върху който се
изпраща \code{Close}; иначе всяко изместване би оставяло по едно заявление на сървъра. Ако
сървърът отговори с код 26000, тоест не познава името на кешираното заявление, записът се
изтрива и потокът се изпълнява повторно веднъж. Това се случва след \code{DISCARD ALL} или
след смяна на сървъра зад пулър и извикващият не бива да го вижда.

\textbf{Удостоверяване.} За PostgreSQL са реализирани доверен вход, парола в чист вид, MD5 и
SCRAM-SHA-256~\cite{rfc7677}. Последният не е разширение, а необходимост: PostgreSQL 14 и
следващите го използват по подразбиране, така че без него клиентът не може да се свърже със
сървър, който не е нарочно отслабен. Обменът е проверен срещу публикувания в RFC 7677 пример:
при фиксиран клиентски случаен низ реализацията възпроизвежда буква по буква и клиентското
доказателство, и подписа на сървъра. За MySQL е реализиран \code{mysql\_native\_password},
включително отговорът на искане за смяна на приставката; \code{caching\_sha2\_password} не е
реализирана и това се съобщава с ясна грешка, а не с мълчалив отказ.

\textbf{Обработка на грешките.} Разграничението е по типове. \code{protocol\_error} означава
байтов поток извън спецификацията и е неотстранима: позицията в потока е загубена.
\code{io\_error} е отказ на транспорта. \code{server\_exception} носи полетата на грешката и е
единствената, след която връзката се използва пак. Връзка с протоколна грешка или отказ на
транспорта се отбелязва като негодна чрез флага \code{poisoned\_}, а пулът проверява
\code{is\_usable()} при връщане и при заемане; тестът
\code{pool\_discards\_a\_connection\_the\_server\_closed\_underneath\_it} убива връзката,
докато тя е заета. За очаквано отсъствие се използва \code{std::optional}, не изключение.
